\begin{center}
    {\LARGE \textbf{2021无机化学与物理化学}} \\[2ex]
    {\large 时量180分钟 \quad 满分150分}
\end{center}

\vspace{0.5cm}
\noindent\textbf{注意事项:}
\begin{enumerate}[label=\arabic*., parsep=0pt, itemsep=0pt]
    \item 答题前,考生先将自己的姓名、学号、考场号、座位号填写在答题卡上,将条形码准确粘贴在考生信息条形码粘贴区。
    \item 考试结束后,将本试卷与答题纸一并收回。
    \item 回答各个问题时,将答案写在答题纸上,写在本试卷上无效。
\end{enumerate}

\vspace{1cm}
\begin{center}
    {\Large \textbf{无机化学部分(共75分)}}
\end{center}

\noindent\textbf{一、选择题(20 pts.)}

\begin{enumerate}[label=\arabic*., itemsep=1.5ex]
    \item 

\end{enumerate}

\vspace{1em}
\noindent\textbf{二、简要回答下列问题(45 pts.)}

\begin{enumerate}[label=\arabic*., itemsep=1.5ex]
    \item 解释实验现象并写出相关的化学反应方程式 (15 pts.)
    \begin{enumerate}[label=(\arabic*), noitemsep, topsep=0pt, partopsep=0pt]
        \item 向 $\ce{CuCl}$ 中滴加氨水，白色沉淀溶解成无色溶液，后变深蓝色。
        \item 少量硫代硫酸钠溶液和硝酸银溶液反应，先生成白色沉淀，沉淀逐渐变黄变棕最后变黑。
        \item $\ce{KI}$ 与 $\ce{NaNO2}$ 混合溶液加入稀硫酸后溶液立即变黄，$\ce{KI}$ 与 $\ce{NaNO3}$ 混合溶液加入稀硫酸后溶液不变色。
    \end{enumerate}

    \item 为什么与过渡金属配位能力\ce{NH3}<\ce{PH3}，而与H+配位能力\ce{NH3}>\ce{PH3}。(5 pts.)

    \item 解释\ce{[NiCl4]^{2-}}是四边形结构、但\ce{[CuCl4]^{2-}}是平面四边形结构。(5 pts.)

    \item 用五种方法制备\ce{NaNO2}和\ce{NaNO3}给出相关实验现象和相应的化学反应方程式。(5 pts.)

    \item 。(5 pts.)

    \item 以铬铁矿原料制备重铬酸钾。(5 pts.)

    \item 试用 3 种方法鉴别 $\ce{MnO2}$ 和 $\ce{PbO2}$，给出相关实验现象和相应的化学反应方程式。(5 pts.)
\end{enumerate}

\vspace{1em}
\noindent\textbf{三、推断题(10 pts.)}

红色固体 A 溶于热的浓盐酸放出

\vspace{1cm}
\begin{center}
    {\Large \textbf{物理化学部分(共75分)}}
\end{center}

\vspace{0.5cm}
\noindent\textbf{四、选择题(24 pts. total, 2 pts. each)}

\begin{enumerate}[label=\arabic*., itemsep=1.5ex]
    \item NaCl稀溶液的摩尔电导率 $\Lambda_m$ 与 $\ce{Na+}$、$\ce{Cl-}$ 离子的淌度 ($U$) 之间的关系为:
    \begin{tasks}(2)
        \task $\Lambda_m = U_+ + U_-$
        \task $\Lambda_m = U_+/F + U_-/F$
        \task $\Lambda_m = U_+F + U_-F$
        \task $\Lambda_m = 2(U_+ + U_-)$
    \end{tasks}

    \item 等温等压下，在A和B组成的均相体系中，若A的偏摩尔体积随浓度的改变而增加，则B的偏摩尔体积将:
    \begin{tasks}(4)
        \task 增加
        \task 减小
        \task 不变
        \task 不一定
    \end{tasks}

    \item 在25℃时，若要使电池 $\ce{Pb(Hg)}(a_1) \mid \ce{Pb(NO3)2}(aq) \mid \ce{Pb(Hg)}(a_2)$ 的电池电动势 $E$ 为正值，则Pb在汞齐中的活度:
    \begin{tasks}(2)
        \task 一定是 $a_1 > a_2$
        \task 一定是 $a_1 = a_2$
        \task 一定是 $a_1 < a_2$
        \task $a_1$ 和 $a_2$ 都可以任意取值
    \end{tasks}

    \item 在标准压力 $p^\theta$ 和268.15 K下，冰变为水，体系的熵变 $\Delta S_{\text{体}}$ 应为:
    \begin{tasks}(4)
        \task 大于零
        \task 小于零
        \task 等于零
        \task 无法确定
    \end{tasks}

    \item 在HAc解离常数测定的实验中，直接测定的物理量是不同浓度的HAc溶液的:
    \begin{tasks}(4)
        \task 电导率
        \task 电阻
        \task 摩尔电导率
        \task 解离度
    \end{tasks}

    \item 恒压下，无相变的单组分封闭体系的焓值随温度的升高而:
    \begin{tasks}(4)
        \task 增加
        \task 减少
        \task 不变
        \task 不一定
    \end{tasks}

    \item 在等温和存在非体积功的情况下，某气体由始态经一微小变化到达终态，其吉布斯自由能的变化公式是:
    \begin{tasks}(2)
        \task $dG = pdV + Vdp + \delta W_f$
        \task $dG = Vdp - \delta W_f$
        \task $dG = pdV + Vdp + \delta W_f$
        \task $dG = -pdV + Vdp - \delta W_f$
    \end{tasks}

    \item 今有298 K, $p^\ominus$ 的 $\ce{N2}$ 气（状态I）和323 K, $p^\ominus$ 的 $\ce{N2}$ 气（状态II）各一瓶，问哪瓶 $\ce{N2}$ 气的化学势大?
    \begin{tasks}(4)
        \task $\mu(\text{I}) > \mu(\text{II})$
        \task $\mu(\text{I}) < \mu(\text{II})$
        \task $\mu(\text{I}) = \mu(\text{II})$
        \task 不可比较
    \end{tasks}

    \item 下列不同浓度的NaCl溶液中（浓度单位 $\text{mol·dm}^{-3}$），哪个溶液的电导率最大?
    \begin{tasks}(4)
        \task 0.001
        \task 0.01
        \task 0.1
        \task 1.0
    \end{tasks}

    \item 下列电解质水溶液中摩尔电导率最大的是:
    \begin{tasks}(2)
        \task $0.001\ \text{mol·kg}^{-1}\ \ce{HAc}$
        \task $0.001\ \text{mol·kg}^{-1}\ \ce{KCl}$
        \task $0.001\ \text{mol·kg}^{-1}\ \ce{KOH}$
        \task $0.001\ \text{mol·kg}^{-1}\ \ce{HCl}$
    \end{tasks}

    \item 下列物理量除哪一个外，均与粒子的阿伏加德罗数有关:
    \begin{tasks}(2)
        \task 法拉第
        \task 标准状态下 $22.4\ \text{dm}^3$ 气体
        \task 摩尔
        \task 库仑
    \end{tasks}

    \item 质量摩尔浓度为 $m$ 的 $\ce{Na3PO4}$ 溶液，平均活度系数为 $\gamma_\pm$，则电解质的活度为:
    \begin{tasks}(2)
        \task $a_B = 4(m/m^\ominus)^4(\gamma_\pm)^4$
        \task $a_B = 4(m/m^\ominus)^4(\gamma_\pm)^4$
        \task $a_B = 27(m/m^\ominus)^4(\gamma_\pm)^4$
        \task $a_B = 27(m/m^\ominus)(\gamma_\pm)^4$
    \end{tasks}
\end{enumerate}

\vspace{0.5cm}
\noindent\textbf{五、计算题(15 pts.)}

1 mol $\ce{H2(g)}$ 从100 K, $4.1\ \text{dm}^3$ 加热到600 K, $49.2\ \text{dm}^3$。若此过程是将气体置于750 K炉中，让其反抗恒定外压 $p^\theta$ 以不可逆方式进行，计算隔离体系的熵变。
\\
已知 $\ce{H2(g)}$ 的摩尔定容热容为:
\[
C_{V,\text{m}} = [20.753 - 0.8368\times10^{-3}(T/\text{K}) + 20.117\times10^{-7}(T/\text{K})^2]\ \text{J·K}^{-1}·\text{mol}^{-1}
\]

\vspace{0.5cm}
\noindent\textbf{六、计算题(10 pts.)}

在413.15 K时，纯 $\ce{C6H5Cl}$ 和纯 $\ce{C6H5Br}$ 的蒸气压分别为125.238 kPa和66.104 kPa。假定两液体组成理想溶液。若有一混合液，在413.15 K, 101.325 kPa下沸腾，试求该溶液的组成，以及在此情况下液面上蒸气的组成。

\vspace{0.5cm}
\noindent\textbf{七、计算题(13 pts.)}

今在300 K时，研究反应 $\ce{A2 + 2B -> 2C + 2D}$，经实验研究得到反应速率方程为 $r = k[\ce{A2}]^{1/2}[\ce{B}]$。
\begin{enumerate}
    \item 320 K时，测得反应之半寿期为300 K时之1/10（以化学计量比投料），求实验活化能。
    \item 请推测可能的反应历程。
\end{enumerate}

\vspace{0.5cm}
\noindent\textbf{八、计算题(12 pts.)}

电池 $\ce{Pt} \mid \ce{H2(g)} \mid \ce{HCl}(0.01\ \text{mol·kg}^{-1}) \mid \ce{AgCl} \mid \ce{Ag(s)}$ 的电动势
\[
E/\text{V} = -0.096 + 1.90\times10^{-3}(T/\text{K}) - 3.041\times10^{-6}(T/\text{K})^2
\]
计算298 K时电池反应的 $\Delta_r G_m$、$\Delta_r S_m$ 和 $\Delta C_{p,\text{m}}$。